\documentclass[11pt]{article}
\usepackage[margin=1in]{geometry}
\usepackage{amsmath,amssymb}
\usepackage{booktabs}
\usepackage{hyperref}
\usepackage{xcolor}
\usepackage{longtable}

\title{Independent Replication Report:\\
\textit{Impact of DNA Repair Kinetics and Dose Rate on RBE Predictions in the UNIVERSE}\\
Liew et al., Int.\ J.\ Mol.\ Sci.\ \textbf{23}, 6268 (2022)}
\author{OpenClaw replication subagent (LUCID set)}
\date{2026-05-30; backfill 2026-07-06}

\begin{document}
\maketitle

\begin{abstract}
We independently reimplement the \emph{photon-side} DNA-repair-kinetics sub-model
of UNIVERSE from the paper text alone (Sec.~5.2 and Eq.~5) and reproduce the
$R_{TD50}$ dose-rate correction for rat spinal cord (RSC) across
14 dose-rate $\times$ fractionation conditions (proton- and helium-SOBP delivery
rates, 1 and 2 fractions), obtaining mean absolute deviations of
0.31\%--1.26\% (overall MAD 0.83\%) against the values printed in the paper's
Table~3 (4th column) and Fig.~4 (left panel). The Table~2 saturation-gain
trend versus dose is reproduced quantitatively at the photon-only
(LET$\to$0) limit for 2, 6, and 12~Gy; at 24~Gy our value is $\sim$7~pp below
the paper's 2~keV/$\mu$m column, consistent with residual proton-track-structure
contribution at low LET. The full proton/helium SOBP RBE curves
(Figs.~1, 2, 4 middle/right, 5) are \textbf{not} reproducible from the paper
alone: they require the closed-source FLUKA HIT beamline geometry, per-spill
delivery logs, and the Kiefer--Chatterjee RDD normalization constants and
Friedrich-2015 saturation formula that the paper cites but does not print in
runnable form. We therefore file this replication as \textbf{PARTIAL}: the new
physics introduced by the paper (photon-side repair kinetics) is replicated
quantitatively; the downstream ion-beam benchmark is not.
\end{abstract}

\section{Verdict}

\textbf{PARTIAL REPLICATION.} The paper's headline mechanism ---
adding time-resolved DNA-DSB repair kinetics to UNIVERSE and thereby
generating a small but well-defined dose-rate correction $R_{TD50}$ to the
photon reference --- is quantitatively replicated ($<$1.3\% MAD across
14 conditions). The downstream ion-beam RBE benchmark (Figs.~1--2, 4
middle/right, 5) is not, and cannot be from the paper alone.

\emph{Headline exercised?} Partially. The new-physics core (photon-side
repair kinetics + $R_{TD50}$) is exercised in full; the ion-side
Kiefer--Chatterjee track-structure sub-model and the FLUKA SOBP delivery
that produce the paper's showcase RBE-vs-dose-rate curves are not.

\section{Reproduced quantitatively}

\begin{itemize}
  \item Photon-side repair-kinetics survival model (Sec.~5.2, Eq.~5), reimplemented
        from scratch.
  \item $R_{TD50}$(dose-rate) for RSC, Fig.~4 left / Table~3 col.~4:
        MAD $0.31\%$--$1.26\%$ across 14 conditions.
  \item Table~2 saturation gain vs.\ dose, LET$\to$0 (photon-only) limit,
        at 2, 6, 12~Gy.
  \item LQ sanity check on DU145 at 2~Gy, 2~Gy/min: $S=0.640$ vs.\ LQ 0.622
        (2.8\% agreement).
\end{itemize}

\section{Not reproduced (and why)}

\begin{itemize}
  \item \textbf{Ion track-structure sub-model} (Kiefer--Chatterjee RDD,
        Eqs.~6--10). Paper prints the functional form; the normalization
        constant $K_p$ depends on a nucleus radius and saturation-energy
        choice not numerically printed for the runs in this paper.
  \item \textbf{Friedrich-2015 LET-dependent DSB-yield boost}, cited
        (ref.~[62]) but not written out.
  \item \textbf{FLUKA HIT scanned-SOBP beamline geometry} and per-spill
        timing logs (29/31 energy slices, 3.5~s spill-to-spill).
        Proprietary to Heidelberg Ion-Beam Therapy Center. Not on any
        public repository.
  \item \textbf{Full proton/helium RBE-vs-dose-rate curves} (Figs.~1, 2,
        4 middle/right, 5) --- ion side is absent so these are not attempted.
  \item \textbf{Cell-line-specific parameter fits}: we used the Table~1
        DU145 and RSC parameter values verbatim; we did \emph{not}
        independently refit any repair rate constants to primary data.
        A genuine reimplementation of the paper's parameter identification
        (chi-square minimization against Karger 2003/2006 + Saager 2018 +
        Hintz 2022 TD50 tables) would be a substantially larger project.
  \item \textbf{No code, no MC files, no fitted-parameter posteriors} were
        released with the paper (``Data Availability Statement: Not applicable'').
\end{itemize}

\section{Honest critique}

\textbf{What was actually exercised.}
The MC engine we built reproduces the photon-side per-cell survival curve
under time-resolved DSB dynamics with (i) Poisson-sampled induction,
(ii) uniform domain distribution across $N_{\text{dom}}=3200$ ``giant-loop''
domains, (iii) same-domain collision reclassification (isolated $\to$ complex
DSB) driven by the printed $\alpha_{\text{DSB}}=30$/(Gy$\cdot$cell), and
(iv) exponentially-distributed repair lifetimes governed by $T_{iDSB}^{1/2}$
and $T_{cDSB}^{1/2}$. Averaging Eq.~5 across 600 MC iterations recovers
$R_{TD50}$ within the model's own Monte-Carlo noise floor
($\sim 0.8\%$--$1.0\%$ one-sigma), which coincides with the observed 0.83\%
overall MAD. That is genuine agreement, not overfitting: we did not tune any
free parameter --- Table~1 values are used verbatim and $N_{\text{dom}}=3200$
is the canonical value from the Friedrich--Scholz lineage that this paper
inherits.

\textbf{What was NOT exercised.} Everything downstream of the photon-side
sub-model. The paper's flagship figure (Fig.~4 middle/right: proton/helium
RBE-weighted dose-rate correction against the Saager 2018 / Hintz 2022 in-vivo
TD50 data) requires the FLUKA-simulated SOBP LET spectrum and delivery
timing, then the ion track-structure sub-model on top. We touched \emph{none}
of that. The claim that ``we replicated the paper'' would be dishonest if
extended past the photon reference channel.

\textbf{Ambiguities in the paper text that biased the result.}
(a) $N_{\text{dom}}$ is not printed; we used the canonical LEM-lineage value.
(b) The paper's Sec.~5.2 does not specify whether an already-classified
complex-DSB domain remains complex after one of its two breaks repairs, or
gets reclassified back to isolated. We took the strict (no-downgrade) reading;
a reclassification policy would perturb $K_{cDSB}$-dominated terms by an
estimated $<1\%$ on $R_{TD50}$.
(c) The reference photon dose rate used to \emph{define} $R_{TD50}$ shifts
between figures (2~Gy/min in Figs.~1--3, Fig.~4 right; 3.75~Gy/min in Fig.~4
left / Table~3). We adopted 3.75~Gy/min for the RSC benchmark because that is
what the numerical $R_{TD50}$ table is normalized to.

\textbf{Overall confidence.} High on the photon-side sub-model (quantitative
agreement, no free parameters). Zero on the ion-side benchmark (not
attempted). ``PARTIAL'' is the honest verdict.

\section{Reproducibility scoring (LUCID rubric)}

\begin{tabular}{lr}
\toprule
Category & Score \\
\midrule
Documentation completeness            & 7 / 10 \\
Independent reimplementation (photon-only sub-model)  & 8 / 10 \\
Independent reimplementation (full proton/helium SOBP) & 3 / 10 \\
Quantitative agreement (achievable subset) & 9 / 10 \\
Open-data score                       & 0 / 10 \\
\bottomrule
\end{tabular}

\section{Open questions}

See \texttt{open\_questions.json} and \texttt{open\_questions\_section.tex}.

\end{document}
